%%%%%%%%%%%%%%%%%%%% author.tex %%%%%%%%%%%%%%%%%%%%%%%%%%%%%%%%%%%
%
% sample root file for your "contribution" to a proceedings volume
%
% Use this file as a template for your own input.
%
%%%%%%%%%%%%%%%% Springer %%%%%%%%%%%%%%%%%%%%%%%%%%%%%%%%%%


\documentclass{svproc}
%
% RECOMMENDED %%%%%%%%%%%%%%%%%%%%%%%%%%%%%%%%%%%%%%%%%%%%%%%%%%%
%

\usepackage{graphicx,amsfonts,amssymb,amsmath,newlfont}
\usepackage{epsfig,url}
\usepackage{enumerate,enumitem}
\usepackage[colorlinks=true,linkcolor=red,citecolor=blue]{hyperref}
\usepackage[dvipsnames]{xcolor}
\usepackage{color}

%Clever ref
\usepackage[noabbrev,capitalize]{cleveref}

\usepackage[all,2cell]{xy} \UseAllTwocells \SilentMatrices

\usepackage{pstricks,pst-node,pst-tree}


% MATH -----------------------------------------------------------
\newcommand{\norm}[1]{\left\Vert#1\right\Vert}
\newcommand{\abs}[1]{\left\vert#1\right\vert}
\newcommand{\set}[1]{\left\{#1\right\}}

\newcommand{\To}{\longrightarrow}
\newcommand*{\Longhookrightarrow}{\ensuremath{\lhook\joinrel\relbar\joinrel\rightarrow}}
\newcommand{\Z}{\mathbb Z}
\newcommand{\Q}{\mathbb Q}
\newcommand{\C}{\mathbb C}
\newcommand{\Ok}{\mathcal O}
\newcommand{\ai}{\mathfrak{a}}
\newcommand{\bi}{\mathfrak{b}}
\newcommand{\R}{\mathbb R}
\newcommand{\N}{\mathbb N}
\newcommand{\AM}{A}
\newcommand{\xx}{\mathsf{x}}
\newcommand{\eqv}{\mathrm{ev}}
\font \rus= wncyr10
\newcommand{\sha}{\, \hbox{\rus x} \,}

\newcommand{\ari}{\mathrm{ar}} %arity
\newcommand{\obj}{\mathrm{Ob}} %object

\newcommand{\GC}{\mathcal{GC}}
\newcommand{\q}{/\!/}

\newcommand{\tr}{\mathrm{tr}}
\newcommand{\id}{\mathrm{id}}

\newcommand{\can}{\mathrm{can}}

\newcommand{\mm}{\mathfrak{m}}

\newcommand{\GL}{\mathrm{GL}}
\newcommand{\LP}{L}
\newcommand{\FL}{F\!L}
\newcommand{\mc}{\mu}

\newcommand{\0}{\color{blue}{\mathsf{0}}}

%%%%Macros Vincent
\newcommand{\eqdef}{\mbox{\,\raisebox{0.2ex}{\scriptsize\ensuremath{\mathrm:}}\ensuremath{=}\,}} % :=
\newcommand{\defeq}{\mbox{~\ensuremath{=}\raisebox{0.2ex}{\scriptsize\ensuremath{\mathrm:}} }} % =:

%%%%Macros PL
\newcommand{\Alt}{ \mid\!\!\mid  } 
\newcommand{\inc}{\subseteq}
\newcommand{\incs}{\subsetneq}
\newcommand{\union}{\cup}
\newcommand{\Union}{\bigcup}	
\newcommand{\comp}{\circ}
\newcommand{\setc}[2]{\set{#1 \mid #2}}
\newcommand \seq[2]{\shortstack{$#1$ \\ \mbox{}\\
                    \mbox{}\hrulefill\mbox{}\\ \mbox{}\\ $#2$}}
\newcommand{\cat}[1]{{\mathbb #1}}
\newcommand{\dl}{[\![} 			
\newcommand{\dr}{]\!]} 
\newcommand{\hyper}[1]{{\mathbb #1}}	
\newcommand{\restrH}[2]{\hyper{#1}\backslash #2}

%Operades

\def\calO{\mathcal{O}}
\newcommand{\KK}{\mathbb{K}}
\newcommand{\opd}[1]{\mathcal{#1}}

%Antischrieck
\newcommand{\as}{{\scriptstyle \text{\rm !`}}}

%Definitions
\definecolor{darkblue}{rgb}{0,0,0.7} % darkblue color
\newcommand{\darkblue}{\color{darkblue}} % darkblue command
\newcommand{\defn}[1]{{\darkblue \emph{#1}}}

%Commentaires 

\newcommand{\Guillaume}[1]{\textcolor{magenta}{\underline{Guillaume}: #1}}
\newcommand{\correction}[1]{\textcolor{red}{#1}}
\newcommand{\correctionbis}[1]{\textcolor{magenta}{#1}}

%Source and sink
\newcommand{\so}{\mathrm{sc}} 
\newcommand{\sk}{\mathrm{sk}} 
\newcommand{\op}{\mathrm{op}}

%PL
\newcommand{\occ}[2]{#1/#2}
\newcommand{\recrestr}[2]{#1_{\cap #2}}
\newcommand{\xyz}[3]{#1\stackrel{#2}{\rightsquigarrow}#3}
\newcommand{\pre}[2]{<_{#1,#2}}

\newcommand{\PL}[1]{{\color{red}{#1}}}

\newcommand{\calP}{\mathcal{P}}
\newcommand{\calH}{\mathcal{H}}
\newcommand{\calT}{\mathcal{T}}
\newcommand{\calU}{\mathcal{U}}

%Operators
\newcommand{\Sat}{\mathrm{Sat}}
\newcommand{\supp}{\mathrm{supp}}
\newcommand{\Ter}{\mathrm{Ter}}
\newcommand{\outsort}{\mathrm{out}}
\newcommand{\insort}{\mathrm{in}}
\newcommand{\var}{\mathrm{var}}

%Dessins

\usepackage{tikz}
\usepackage{tikz-cd}
\usepackage{pgfplots}
\usepackage{pgfplotstable}
\tikzset{math3d/.style=
    {x= {(-0.353cm,-0.353cm)}, z={(0cm,1cm)},y={(1cm,0cm)}}}
\tikzset{JLL3d/.style=
    {x= {(0.4cm,-0.2cm)}, z={(0cm,1cm)},y={(-1cm,0cm)}}}
\usetikzlibrary{calc}
\usetikzlibrary{shapes,shapes.geometric,fit,positioning,calc,matrix}
\tikzset{
  optree/.style={scale=.5,thick,grow'=up,level distance=10mm,inner sep=1pt},
  comp/.style={draw=none,circle,fill,line width=0,inner sep=0pt},
  dot/.style={draw,circle,fill,inner sep=0pt,minimum width=3pt},
  circ/.style={draw,circle,inner sep=1pt,minimum width=4mm},
  emptycirc/.style={draw,circle,inner sep=1pt,minimum width=2mm},
  root/.style={level distance=10mm,inner sep=1pt},
  leaf/.style={draw=none,circle,fill,line width=0,inner sep=0pt},
  nodot/.style={draw,circle,inner sep=1pt},
}

\pgfplotsset{compat=1.12}


% to typeset URLs, URIs, and DOIs
\usepackage{url}
\def\UrlFont{\rmfamily}

\begin{document}


\mainmatter              % start of a contribution
%
\title{Term Rewriting on Nestohedra}
%
\titlerunning{Term Rewriting on Nestohedra}  % abbreviated title (for running head)
%                                     also used for the TOC unless
%                                     \toctitle is used
%
\author{Pierre-Louis Curien\inst{1} \and Guillaume Laplante-Anfossi\inst{2}}
%
\authorrunning{Pierre-Louis Curien and Guillaume Laplante-Anfossi} % abbreviated author list (for running head)
%
%%%% list of authors for the TOC (use if author list has to be modified)
\tocauthor{Pierre-Louis Curien and Guillaume Laplante-Anfossi}
%
\institute{IRIF, Universit\'e Paris Diderot and $Picube$ team, Inria, France,\\
\email{curien@irif.fr},\\ ORCiD: 0009-0000-1226-9574
\and
Centre for Quantum Mathematics, IMADA, Syddansk Universitet, Odense, Denmark,\\
\email{glaplanteanfossi@imada.sdu.dk},\\ ORCiD: 0000-0002-7383-4474}

\maketitle              % typeset the title of the contribution

\begin{abstract}
  We define term rewriting systems on the vertices and faces of nestohedra, and show that the former are  
  confluent and terminating. 
  While the associated posets on vertices generalize Barnard--McConville's flip order for graph-associahedra, the preorders on faces generalize the facial weak order for permutahedra and the generalized Tamari order for associahedra. 
  Moreover, we define and study contextual families of nestohedra, whose local confluence diagrams satisfy a certain uniformity condition. 
  Among them are associahedra and operahedra, whose associated proofs of confluence for their rewriting systems reproduce proofs of categorical coherence theorems for monoidal categories and categorified operads.
% We would like to encourage you to list your keywords within
% the abstract section using the \keywords{...} command.
\keywords{term rewriting, nestohedra, facial weak order, generalized Tamari order, categorical coherence}
\end{abstract}

\input{sec/introduction2}
\input{sec/nestohedra}
\input{sec/2faces}
\input{sec/rewriting2}
\input{sec/contextual}

%
% ---- Bibliography ----
%
\bibliographystyle{amsalpha}

\bibliography{Coherence2}

\end{document}
